% Self-contained report. Compile with pdfLaTeX in Overleaf (twice for contents).
\documentclass[12pt,a4paper]{article}
\usepackage[T1]{fontenc}
\usepackage[utf8]{inputenc}
\usepackage[english]{babel}
\usepackage[margin=2.5cm]{geometry}
\usepackage{xcolor}
\usepackage{listings}
\usepackage{booktabs}
\usepackage{microtype}
\usepackage[hidelinks]{hyperref}
\definecolor{listinggray}{RGB}{247,247,247}
\lstset{basicstyle=\ttfamily\footnotesize,backgroundcolor=\color{listinggray},
 frame=single,rulecolor=\color{gray},breaklines=true,breakatwhitespace=false,
 columns=fullflexible,keepspaces=true,showstringspaces=false,tabsize=4,
 captionpos=b,aboveskip=10pt,belowskip=10pt}
\setlength{\emergencystretch}{3em}
\begin{document}
\begin{titlepage}
\centering
\textbf{Ministry of Education, Culture and Research of the Republic of Moldova}\par\medskip
\textbf{Technical University of Moldova}\par\medskip
\textbf{Faculty of Computers, Informatics and Microelectronics}\par\medskip
\textbf{Department of Software Engineering and Automatic}\par
\vspace{2.5cm}
\textbf{ANDREI ELCISCEV}\par
\textbf{Group FAF-243}\par
\vspace{2cm}
{\Large\textbf{LABORATORY WORK No. 17}\par}
\medskip
{\Large\textbf{THREADS SHARE MEMORY, AND RACE}\par}
\vspace{2cm}
\begin{flushright}
\textbf{Checked by:}\\
Rei\c{t}man Patricia
\end{flushright}
\vfill
\textbf{Chi\c{s}in\u{a}u, 2026}
\end{titlepage}
\pagenumbering{arabic}
\tableofcontents
\newpage
\section{Introduction}
This laboratory work examines how POSIX threads share a process's memory and how
unsynchronized updates can be lost. Each worker increments one global counter a
million times. The experiment first uses two workers, then increases their number
to four.

The work ran in \texttt{/home/user/os-lab17} in the Linux VM. The compiler used
\texttt{-Wall -Wextra -pthread} at its default optimization level. The listings
contain actual commands and output. The original two-thread version is retained
as \texttt{fire\_two.c}, and \texttt{fire.c} contains the four-thread extension.

\begin{lstlisting}[caption={Working directory, kernel, available CPUs, and compiler}]
$ pwd
/home/user/os-lab17
$ uname -sr
Linux 6.14.0-37-generic
$ nproc
3
$ gcc -dumpfullversion
13.3.0
\end{lstlisting}

\section{Part 1: Shared and Racing}
The main thread creates two workers and joins both before printing the final
counter. Joining waits for completion, while the workers access the counter
without synchronization.

\begin{lstlisting}[caption={Three actual trials with two workers}]
$ gcc -Wall -Wextra -pthread -o fire_two fire_two.c
$ ./fire_two
thread 2 sees counter = 0
thread 1 sees counter = 0
expected: 2000000
counter:  1207310   (lost updates: the race)
$ ./fire_two
thread 1 sees counter = 0
thread 2 sees counter = 0
expected: 2000000
counter:  1321397   (lost updates: the race)
$ ./fire_two
thread 1 sees counter = 0
thread 2 sees counter = 195948
expected: 2000000
counter:  1821131   (lost updates: the race)
\end{lstlisting}

\subsection{Observe: The Same Variable}
In the third trial, thread 2 printed a starting value of 195948. Thread 1 had
already started incrementing the same global counter. Both threads operate in
one address space. In the other two trials, both initial reads were zero because
the relative timing differed.

\subsection{Observe: Lost Increments}
An increment reads the value, adds one, and stores the result. If both threads
read 10 before either stores, both can compute 11 and store 11. Two attempted
increments then produce a net increase of one. Scheduling and parallel execution
change how these operations interleave.

The conflicting unsynchronized accesses constitute a C data race with undefined
behavior. The measured outcomes illustrate lost updates. A specific count, a
different count in every run, or a count below the expected total in every
execution is not guaranteed by the language.

\section{Part 2: Four Workers}
Two additional workers use IDs 3 and 4. The main thread joins all four and compares
the result with \(4\times1{,}000{,}000=4{,}000{,}000\).

\begin{lstlisting}[caption={Three actual trials with four workers}]
$ gcc -Wall -Wextra -pthread -o fire fire.c
$ ./fire
thread 1 sees counter = 0
thread 2 sees counter = 87834
thread 3 sees counter = 1010351
thread 4 sees counter = 1131088
expected: 4000000
counter:  2417179   (lost updates: the race)
$ ./fire
thread 1 sees counter = 0
thread 2 sees counter = 67579
thread 3 sees counter = 283743
thread 4 sees counter = 337420
expected: 4000000
counter:  2087047   (lost updates: the race)
$ ./fire
thread 2 sees counter = 0
thread 1 sees counter = 0
thread 3 sees counter = 48173
thread 4 sees counter = 235861
expected: 4000000
counter:  1629452   (lost updates: the race)
\end{lstlisting}

\subsection{Comparing Lost Updates}
Each loss count below is the expected total minus the observed total. Each row
describes a separate execution.
\begin{center}
\begin{tabular}{rrrrr}
\toprule
Workers & Trial & Expected & Observed & Lost updates\\
\midrule
2 & 1 & 2000000 & 1207310 & 792690\\
2 & 2 & 2000000 & 1321397 & 678603\\
2 & 3 & 2000000 & 1821131 & 178869\\
4 & 1 & 4000000 & 2417179 & 1582821\\
4 & 2 & 4000000 & 2087047 & 1912953\\
4 & 3 & 4000000 & 1629452 & 2370548\\
\bottomrule
\end{tabular}
\end{center}
The mean absolute loss rose from 550054 to 1955441 updates. The mean lost fraction rose from 27.50\% to 48.89\%.

The VM exposes three CPUs. The extra workers increased loss in these measurements.
The number lost depends on timing, so individual executions can differ.

\subsection{Observe: One Idea for a Safe Increment}
Protect the entire read--add--write operation with one shared mutex so only one
worker can increment the counter at a time. The handout asks for a proposal,
so the experimental programs retain the unprotected increment.

\section{Conclusion}
The initial counter reads and the final counts demonstrated shared memory and
lost updates in the recorded runs. Increasing from two to four workers increased
both the mean number and mean fraction of lost updates in this VM. Mutual
exclusion around the complete increment would serialize that critical section.

\appendix
\section{Submitted Source Code}

\begin{lstlisting}[caption={fire\_two.c},language=C]
/* fire.c - two threads of one process share the same memory.
 * Both add to the same counter, a million times each, WITHOUT a lock.
 * Build: gcc -Wall -Wextra -pthread -o fire fire.c   Run: ./fire
 */
#include <stdio.h>
#include <pthread.h>

#define N 1000000
long counter = 0;              /* one variable, shared by both threads */

void *work(void *arg) {
    long id = (long)arg;
    printf("thread %ld sees counter = %ld\n", id, counter);
    for (int i = 0; i < N; i++)
        counter++;             /* NOT atomic: read, add, write */
    return NULL;
}

int main(void) {
    pthread_t a, b;
    pthread_create(&a, NULL, work, (void *)1);
    pthread_create(&b, NULL, work, (void *)2);
    pthread_join(a, NULL);
    pthread_join(b, NULL);
    printf("expected: %d\n", 2 * N);
    printf("counter:  %ld   (lost updates: the race)\n", counter);
    return 0;
}
\end{lstlisting}

\begin{lstlisting}[caption={fire.c},language=C]
/* fire.c - four threads of one process share the same memory.
 * All add to the same counter, a million times each, WITHOUT a lock.
 * Build: gcc -Wall -Wextra -pthread -o fire fire.c   Run: ./fire
 */
#include <stdio.h>
#include <pthread.h>

#define N 1000000
long counter = 0;              /* one variable, shared by all four threads */

void *work(void *arg) {
    long id = (long)arg;
    printf("thread %ld sees counter = %ld\n", id, counter);
    for (int i = 0; i < N; i++)
        counter++;             /* NOT atomic: read, add, write */
    return NULL;
}

int main(void) {
    pthread_t a, b, c, d;
    pthread_create(&a, NULL, work, (void *)1);
    pthread_create(&b, NULL, work, (void *)2);
    pthread_create(&c, NULL, work, (void *)3);
    pthread_create(&d, NULL, work, (void *)4);
    pthread_join(a, NULL);
    pthread_join(b, NULL);
    pthread_join(c, NULL);
    pthread_join(d, NULL);
    printf("expected: %d\n", 4 * N);
    printf("counter:  %ld   (lost updates: the race)\n", counter);
    return 0;
}
\end{lstlisting}

\end{document}
